\section{Input file formats, in full}
\label{sec:input-files}

Everything the program reads is a plain text, JSON, XML or FITS file that
you can inspect and produce yourself. This section specifies each format
completely: the exact columns, their units, the parsing rules, and a
literal example you can copy. Blank lines are ignored everywhere; lines
beginning with \code{\#} are comments where noted.

\subsection{The template catalogue: \code{data/interpola.db.csv}}
\label{sec:file-catalog}

This comma-separated file is the master list of stellar templates; the
template selector in the GUI is built from it. The first line is a header
and is skipped. Every following line describes one template with exactly
six fields:

\begin{center}
\begin{tabular}{clp{8.6cm}}
\toprule
\# & Field & Meaning \\
\midrule
1 & \code{nrow} & Integer: the number of data rows in the spectrum file.
    Used only as a consistency check at load time (a mismatch larger than
    one row prints a warning, nothing more), so an approximate value is
    acceptable. \\
2 & \code{Vmag} & Float: the \emph{visual magnitude that the spectrum file
    represents}, called $m_{v0}$ throughout the documentation. This is the
    heart of the calibration: the ETC multiplies the file by
    $10^{0.4\,m_{v0}}$ to bring it to the ``visual zero'' scale and then
    scales it to the magnitude you enter in the GUI. For an absolutely
    calibrated file (CALSPEC), $m_{v0}$ is simply the synthetic $V$
    magnitude of the file itself: $0.026$ for Vega, $-26.75$ for the Sun.
    For the BPGS library the files are already normalised to visual zero,
    hence $m_{v0}=0$. If this number is wrong, every prediction for that
    template is wrong by the same factor --- when in doubt, let
    \code{add\_template.py} compute it (Sect.~\ref{sec:tool-template}). \\
3 & \code{B-V} & Float: the Johnson $B{-}V$ colour of the template.
    Informational: it is displayed in the selector and used by the
    template-validation button, but it does not enter the flux
    calibration. \\
4 & \code{name} & The display name shown in the selector. Any text without
    a comma; leading/trailing spaces are stripped. \\
5 & \code{spt} & The spectral type (\code{A0V}, \code{G2V}, \code{M5III},
    \dots). The search box of the selector matches against this field. \\
6 & \code{filename} & Path of the spectrum file, \emph{relative to the
    data directory}. Subdirectories are allowed and encouraged
    (\code{bpgs/bpgs\_92.ascii}, \code{imported/hd12345.fits}). \\
\bottomrule
\end{tabular}
\end{center}

A literal excerpt:

\begin{verbatim}
nrow ,Vmag   ,B-V     ,name       ,spt   ,filename
 8839,  0.026,   0.000, Vega      ,   A0V, bpgs/alpha_lyr_stis_004.ascii
 1467,-26.750,   0.630, Sun       ,   G2V, bpgs/sun_reference_stis_001.ascii
  481,  0.000,   0.310, HD 76483  ,   A3V, bpgs/bpgs_48.ascii
\end{verbatim}

Fields are separated by commas; the spaces you see are padding for human
readability and are stripped on parsing. A line with fewer than six
fields, or with non-numeric values in the numeric fields, is silently
skipped --- a convenient way to comment a template out is therefore to
put a letter in front of its \code{nrow}.

\subsection{Template spectra}
\label{sec:file-templates}

Two formats are accepted, distinguished by the file extension.

\paragraph{Two-column ASCII (any extension except \code{.fits}/\code{.fit}).}
Each line holds a wavelength in \emph{\AA{}ngstr\"om} and a flux density
$F_\lambda$ in \si{erg.s^{-1}.cm^{-2}.\angstrom^{-1}}, separated by
whitespace:

\begin{verbatim}
 3500.0  4.215e-09
 3502.5  4.198e-09
 3505.0  4.187e-09
\end{verbatim}

Parsing rules, all applied automatically at load: lines with fewer than
two numeric tokens are skipped; rows with non-finite, zero or negative
wavelength or flux are dropped; rows are sorted by wavelength; duplicate
wavelengths are removed keeping the first occurrence. What matters
physically is the \emph{shape}: because the ETC rescales the spectrum to
the magnitude you enter (via $m_{v0}$, previous subsection), a template in
arbitrary flux units is perfectly usable as long as the relative
calibration across wavelength is correct.

\paragraph{CALSPEC/STScI FITS tables (\code{.fits}, \code{.fit}).}
The standard format of the STScI calibration archives: a FITS binary
table with a wavelength column named \code{WAVELENGTH} (or \code{WAVE},
\code{LAMBDA}, \code{WL}) and a flux column named \code{FLUX} (or
\code{FLAM}, \code{SURF\_BRIGHT}, \code{SB}), matched case-insensitively.
The wavelength unit is read from the column's \code{TUNIT} keyword and may
be \AA{}ngstr\"om (the CALSPEC convention), nanometres or microns ---
conversion is automatic. Additional columns (\code{STATERROR},
\code{SYSERROR}, \code{FWHM}, \dots) are ignored. This one reader covers
the CALSPEC standard-star archive, the solar-system surface-brightness
atlas, the galactic emission-line atlas, the supernova/kilonova transient
templates and the CLOUDY planetary-nebula grids without modification.

\subsection{The filter list and profiles}
\label{sec:file-filters}

\paragraph{\code{data/filters.list}} names the available filters, one
entry per line; \code{\#} comments allowed. Each \emph{logical} filter
(say \code{Bessell.V}) requires two profile files, \code{Bessell.V\_Vega}
and \code{Bessell.V\_AB}, holding the same transmission curve with the two
photometric calibrations. Entries may be written with or without the
\code{filters/} prefix and with or without the \code{.xml} extension ---
all four combinations resolve. Example:

\begin{verbatim}
filters/Bessell.V_Vega
filters/Bessell.V_AB
filters/SDSS.r_Vega
filters/SDSS.r_AB
\end{verbatim}

\paragraph{Profile files (\code{data/filters/*.xml})} are VOTables in the
SVO Filter Profile Service format. If you download a filter from SVO
(\url{http://svo2.cab.inta-csic.es/theory/fps/}), take the ``VOTable''
link of the desired calibration and drop the file in unchanged. The
parser requires four \code{PARAM} elements --- \code{MagSys} (\code{Vega}
or \code{AB}), \code{ZeroPoint} with \code{ZeroPointUnit} \code{Jy},
\code{DetectorType} (0 = energy counter, 1 = photon counter) and
\code{WavelengthUnit} --- plus the transmission table itself as
\code{TR}/\code{TD} rows of wavelength and transmission. All other
metadata (pivot wavelength, FWHM, effective width, \dots) is read when
present and shown in the filter panel. For a filter that SVO does not
carry --- every amateur RGB or narrowband filter --- build the pair with
\code{make\_filter\_profile.py} from a plain transmission curve
(Sect.~\ref{sec:tool-filter}); the generated files contain everything
listed above.

\paragraph{\code{BLANK.xml}} is the special ``no filter'' profile: unity
transmission from 3000 to 26000\,\AA{}. Because an unfiltered measurement
has no meaningful Vega calibration, it is defined on the AB scale for
both selector positions. Spectroscopy selects it by default.

\subsection{Detector and instrument curves}
\label{sec:file-curves}

Four optional two-column text files describe your instrument beyond the
scalar numbers in the GUI. All four share the same syntax: whitespace- or
comma-separated columns, \code{\#} comments, and a wavelength unit that
you declare in the GUI (\AA{}ngstr\"om, nm or \si{\micro\metre}) rather
than in the file --- so a curve digitized from a datasheet in nanometres
is used as-is by setting the unit selector to \code{nm}.

\begin{description}
\item[QE curve (\code{data/qe.dat} or per-profile).] Wavelength and
  quantum efficiency as a \emph{fraction} (0--1):
\begin{verbatim}
# ASI533MM  wavelength[nm]  QE
 400  0.62
 500  0.80
 600  0.72
 700  0.52
\end{verbatim}
  Values are clipped to [0, 1]; outside the tabulated range the QE is
  taken as zero, which is physically right for a sensor cut-off. For an
  OSC camera in single-channel mode, this file should be the
  \emph{channel's} effective QE: sensor QE multiplied by the colour-dye
  transmission of that channel, both published by the camera makers.
\item[Optics response curve (optional).] Same syntax; the wavelength-%
  dependent throughput of everything between aperture and detector
  \emph{excluding} the QE and the filter (mirror reflectivity, corrector
  and window transmission, \dots). It multiplies the scalar efficiency
  from the GUI. If you have only the scalar, leave this file out.
\item[Slit resolving-power curve (optional).] Two columns: slit width in
  arcsec, measured resolving power $R$. When loaded, the measured curve
  overrides the nominal $R$ for slit widths inside its range; a width
  outside the range is refused rather than extrapolated.
\item[CMOS gain table (optional).] Four columns: gain \emph{setting} (the
  number you dial into the camera driver), conversion gain in
  e$^-$/ADU, read noise in e$^-$ rms, full-well capacity in e$^-$:
\begin{verbatim}
# ASI533MM: setting  e/ADU   RN    FWC
     0      3.64    3.55   50000
   100      1.00    1.45   13700
   200      0.32    1.20    4400
\end{verbatim}
  Selecting a row in the GUI fills the three detector fields together, so
  gain, noise and saturation always describe the same camera state. The
  numbers come from the maker's published curves or from
  PhotonsToPhotos.
\end{description}

\subsection{The atmospheric transmission FITS file}

\code{data/earth\_atmospheric\_transmission.fits}, if present, replaces
the built-in broad-band extinction table with a measured curve. Three
layouts are accepted: a binary table with a wavelength column
(\code{WAVEAIR}, \code{WAVELENGTH}, \code{WAVE}, \code{LAMBDA},
\code{LAM} or \code{WL}) and a direct transmission column
(\code{TRANSMISSION}, \code{TRANS}, \code{THROUGHPUT}) holding fractions;
the same with an \code{EXTINCTION} column in mag/airmass (the Paranal
product), converted internally as $T = 10^{-0.4 k_\lambda}$; or a plain
$2\times N$ image whose wavelength unit is declared by \code{CUNIT1}.
Whatever the source, the file must state \emph{zenith} values: the ETC
raises the curve to the current airmass itself, exactly once. Outside the
file's wavelength range the edge values are extended (an atmosphere does
not become opaque where a file ends).

\subsection{Sites: \code{observatories.json}}

All sites --- the shipped defaults and any you save --- live in a single
JSON list next to the program, \code{observatories.json}:

\begin{verbatim}
[
  { "name": "My backyard", "lat": 45.4064, "lon": 11.8768,
    "elev": 12, "mag_sky": 19.8,
    "timezone": "Europe/Rome", "utc_offset_h": 1 }
]
\end{verbatim}

\code{lat}/\code{lon} are decimal degrees (north and east positive),
\code{elev} metres, \code{mag\_sky} the default value for the fixed-sky
mode, and \code{timezone}/\code{utc\_offset\_h} set the observation time
reference automatically when the site is selected. You rarely edit this
file by hand: \textsf{Save site} writes the entry for you, and
\textsf{Import\ldots} merges another user's JSON into it. (The former
\code{etc\_sites.json} is no longer used; \code{observatories.json} is now
the single site store.)

\subsection{Instrument profiles}
\label{sec:profiles}

\textsf{Save instrument profile} writes a JSON file containing every
telescope/detector/spectroscopy setting of the middle columns, and copies
the currently loaded QE, optics-response and slit-$R$ curves beside it as
\code{<profile>\_qe.dat}, \code{<profile>\_throughput.dat} and
\code{<profile>\_slitres.dat}, referenced by relative name --- the set of
files is portable as a folder. Loading a profile restores everything,
including the curves; the most recently used profile is reloaded
automatically at start-up. This is the mechanism to maintain one
configuration per telescope+camera combination you own.

\subsection{Horizon profile CSV}
\label{sec:file-horizon}

Written automatically by the horizon generator into
\code{data/horizons/horizon\_lat<lat>\_lon<lon>\_r<km>km.csv}; readable by
any plotting tool and by \textsf{Display horizon}. A commented header
records the site, then one row per azimuth:

\begin{verbatim}
# SPETC horizon profile
# latitude_deg: 45.848600
# longitude_deg: 11.569600
# radius_km: 10
# center_elevation_m: 1366.0
azimuth_deg,horizon_elevation_deg
0.0,1.234
1.0,1.198
\end{verbatim}

Azimuth is degrees from North through East (N=0, E=90 --- the same
convention as every other azimuth in the program); the horizon elevation
is the apparent altitude of the terrain in degrees, including the
Earth-curvature drop and the standard $34'$ horizontal refraction, so a
target is genuinely visible when its computed altitude exceeds the
profile value at its azimuth.

\subsection{Batch run descriptions}

The headless mode (Sect.~\ref{sec:tool-batch}) is driven by a JSON file
whose blocks mirror the GUI columns: \code{telescope} (mm and scalar
efficiency), \code{detector} (the \code{Detector} constructor fields),
\code{atmosphere} (airmass, seeing, PSF model, elevation),
\code{sky} (\code{"mode": "fixed\_ab"} with \code{sky\_mag}, or
\code{"mode": "sqm"} with \code{sqm\_v\_mag\_arcsec2}), \code{target}
(template name from the catalogue, magnitude, system, reference filter),
and a \code{photometry} or \code{spectroscopy} block with the
mode-specific settings. The two shipped examples in \code{examples/} are
complete and commented by construction --- copy one and edit the numbers.
